\chapter{Tensors and Shapes}
\label{ch:tensors}

\chapterorient{\Cref{ch:bigidea} sold you a slogan---\emph{a field is a
tensor}---and then leaned on it without ever pinning down what a tensor
\emph{is}. This chapter pays that debt. We say precisely what a tensor is in our
calculus (an immutable, typed, multi-axis array of numbers), introduce
\emph{shapes} as the type that travels with every tensor and is checked at every
function boundary, and write down the handful of primitive operations---%
\code{axpy}, \code{dot}, \code{matvec}, and friends---each with its shape
contract drawn as a picture. Then comes the chapter's real payload: the
\emph{tensor-field lift}, the modeling move that turns a per-degree-of-freedom
loop over an array (the world of Part~II) into a single equation about a whole
field. By the end you will read a shape the way you read a type, and you will see
why ``loop over the entries'' and ``one statement about the field'' are usually,
but not always, the same thing. The careful \emph{algebra} of shapes---how they
combine, broadcast, and promote---is the next chapter, \cref{ch:shapes}; here we
fix the vocabulary and the simplest shapes.}

\section{What a tensor is}
\label{sec:t-whatis}

We have used the word ``tensor'' since \cref{ch:bigidea} as a synonym for ``a
rectangular block of numbers.'' That is the right intuition, but our calculus
asks three specific things of the word, and it is worth stating them plainly
because each one earns its keep later.

\begin{definition}[Tensor]
\label{def:tensor}
A \emph{tensor} is an \emph{immutable}, \emph{typed}, \emph{multi-axis} array of
scalars:
\begin{description}[leftmargin=1.6em, style=nextline, nosep]
\item[immutable] once constructed, its entries never change. An operation that
``modifies'' a tensor in fact returns a \emph{new} tensor and leaves its input
untouched (\cref{ch:bigidea}).
\item[typed] it carries a \emph{shape}---the number, sizes, and meanings of its
axes---and a \emph{scalar type} (real or complex). The shape is part of the
value, not a comment about it.
\item[multi-axis] its entries are addressed by a tuple of integer indices, one
per axis. A length-$N$ vector has one axis; an $m\times n$ matrix has two; the
$[H,W,3]$ field of \cref{ch:bigidea} has three.
\end{description}
\end{definition}

The \emph{number} of axes is the tensor's \emph{rank}: rank-1 for a vector,
rank-2 for a matrix, rank-0 for a single scalar (\cref{sec:t-scalars}). Do not
confuse this with the rank of a matrix in the linear-algebra sense (the dimension
of its column space); throughout this book ``rank'' means ``number of axes''
unless we say otherwise.

\subsection{Reconnecting to Part~II: every field we built is a tensor}
\label{sec:t-reconnect}

This is not a new object. Every quantity Part~II constructed is already a tensor;
we simply had not insisted on the word. Two cases anchor everything that follows.

\paragraph{A discretized field is a \code{Tensor[N]}.}
\Cref{ch:functionspaces} replaced a continuous field---the potential $V$, the
mode field $\vE$---with a finite list of \emph{degrees of freedom} (DoFs): the
coefficients that, multiplied against the basis functions, reconstruct the field
everywhere. If the function space has $N$ degrees of freedom, that list is $N$
numbers in a fixed order. As a tensor it has exactly one axis, of length $N$:

\begin{center}
\lstinline[style=l4style]|Tensor[N]| \quad---\quad a rank-1 tensor, the global
DoF vector.
\end{center}

\noindent This is \emph{the} workhorse shape of the whole solver. The unknown we
solve for, the right-hand side we solve against, every residual and search
direction a Krylov method forms (\cref{ch:matrixfree})---all are
\lstinline[style=l4style]|Tensor[N]|. The single axis $N$ counts degrees of
freedom, and which DoFs they are is the business of the function space; the
tensor only remembers \emph{how many} and \emph{in what order}.
\Cref{fig:t-meshfield} draws the correspondence: the continuous field over the
mesh on the left, and the flat \lstinline[style=l4style]|Tensor[N]| of its DoF
coefficients on the right.

\begin{figure}[t]
\centering
\begin{tikzpicture}[scale=1.0]
  % LEFT: a triangulated mesh with a few highlighted nodes (DoFs)
  \begin{scope}
    \draw[simGray, semithick]
      (-1.5,-1.0) -- (0.0,-1.1) -- (1.4,-0.8) --
      (1.2,0.7) -- (-0.2,1.0) -- (-1.6,0.6) -- cycle;
    \draw[simGray, semithick] (-1.5,-1.0) -- (-0.2,1.0);
    \draw[simGray, semithick] (0.0,-1.1) -- (-0.2,1.0);
    \draw[simGray, semithick] (0.0,-1.1) -- (1.2,0.7);
    \draw[simGray, semithick] (1.4,-0.8) -- (1.2,0.7);
    \draw[simGray, semithick] (-1.6,0.6) -- (-0.2,1.0);
    \draw[simGray, semithick] (-1.5,-1.0) -- (-1.6,0.6);
    % the DoF nodes, numbered
    \foreach \p/\n in {{(-1.5,-1.0)}/0, {(0.0,-1.1)}/1, {(1.4,-0.8)}/2,
                       {(1.2,0.7)}/3, {(-0.2,1.0)}/4, {(-1.6,0.6)}/5}{
      \fill[simRust] \p circle (2.1pt);
      \node[font=\tiny, simRust, anchor=south west] at \p {\n};}
    \node[font=\scriptsize\itshape, simGray] at (0,-1.6) {a field on a mesh};
    \node[font=\scriptsize\itshape, simGray] at (0,-1.95) {$N=6$ DoFs (one per node)};
  \end{scope}
  % collapse arrow
  \node[font=\Large, simBlue] at (2.7,0) {$\rightsquigarrow$};
  \node[font=\scriptsize\itshape, simBlue] at (2.7,0.55) {collect the};
  \node[font=\scriptsize\itshape, simBlue] at (2.7,-0.55) {DoF coefficients};
  % RIGHT: the Tensor[N] column with the 6 entries
  \begin{scope}[shift={(4.9,0)}]
    \fill[simBgBlue, draw=simBlue, thick] (-0.32,-1.2) rectangle (0.32,1.2);
    \foreach \k in {0,1,2,3,4,5}{
      \pgfmathsetmacro\yy{0.9-0.4*\k}
      \draw[simBlue!35] (-0.32,\yy+0.2)--(0.32,\yy+0.2);
      \node[font=\tiny, simRust] at (-0.62,\yy) {\k};
      \node[font=\scriptsize, simInk] at (0.0,\yy) {$c_{\k}$};}
    \node[font=\scriptsize, simRust] at (0.95,0) {$N$};
    \draw[-{Stealth[length=1.6mm]}, simRust] (0.62,1.0)--(0.62,-1.0);
    \node[font=\scriptsize\ttfamily, simInk] at (0.0,-1.6) {Tensor[N]};
    \node[font=\scriptsize\itshape, simGray] at (0.0,-1.95) {rank-1 DoF vector};
  \end{scope}
\end{tikzpicture}
\caption[A field over a mesh is a \code{Tensor[N]}]{The correspondence that
underwrites the whole calculus. \emph{Left:} a field defined over a mesh, sampled
by $N$ degrees of freedom---here one coefficient per node, $N=6$. \emph{Right:}
those coefficients collected, in their fixed order, into a single rank-1 tensor
\lstinline[style=l4style]|Tensor[N]|. The mesh, the basis, and the geometry choose
\emph{which} DoFs and in \emph{what order} (the function-space machinery of
\cref{ch:functionspaces}); the tensor remembers only the ordered list. From
Part~III's height, the field \emph{is} this \lstinline[style=l4style]|Tensor[N]|.}
\label{fig:t-meshfield}
\end{figure}

\paragraph{An element-local quantity is a small multi-axis tensor.}
\Cref{ch:matrixfree} showed that the moment we gather the global vector down to
individual elements, we acquire \emph{richer} shape. The gathered coefficients
form a \lstinline[style=l4style]|Tensor[E, L]| ($E$ elements, $L$ local DoFs
each); the values at quadrature points form a
\lstinline[style=l4style]|Tensor[E, P, C]| ($P$ points, $C$ components); the
geometry factor is a \lstinline[style=l4style]|Tensor[E, P, G]|. These are the
\emph{same kind of object} as the flat DoF vector---immutable, typed, multi-axis
arrays---only with more axes and more specific meanings. The gather/scatter pair
$G,\,G^{\mathsf T}$ of \cref{ch:matrixfree} is precisely the bridge between the
one-axis global vocabulary and the multi-axis element-local one.

\Cref{fig:t-box} draws three representative tensors as labeled boxes of axes---a
picture we will return to whenever a shape needs to be made concrete.

\begin{figure}[t]
\centering
\begin{tikzpicture}[scale=1.0]
  % --- Tensor[N]: a single tall column ---
  \begin{scope}[shift={(-5.0,0)}]
    \fill[simBgBlue, draw=simBlue, thick] (-0.28,-1.1) rectangle (0.28,1.1);
    \foreach \y in {-0.825,-0.55,-0.275,0,0.275,0.55,0.825}{
      \draw[simBlue!40] (-0.28,\y)--(0.28,\y);}
    \node[font=\scriptsize, simRust] at (0.0,1.45) {$N$};
    \draw[-{Stealth[length=1.6mm]}, simRust] (0.55,0.9)--(0.55,-0.9);
    \node[font=\scriptsize\ttfamily, simInk] at (0.0,-1.5) {Tensor[N]};
    \node[font=\scriptsize\itshape, simGray, text width=26mm, align=center]
      at (0.0,-2.05) {global DoF vector\\(rank 1)};
  \end{scope}
  % --- Tensor[m, n]: a small grid ---
  \begin{scope}[shift={(-1.4,0)}]
    \fill[simBgGreen, draw=simGreen, thick] (-0.85,-0.7) rectangle (0.85,0.7);
    \foreach \y in {-0.35,0,0.35}{\draw[simGreen!45] (-0.85,\y)--(0.85,\y);}
    \foreach \x in {-0.51,-0.17,0.17,0.51}{\draw[simGreen!45] (\x,-0.7)--(\x,0.7);}
    \node[font=\scriptsize, simInk] at (-1.15,0.0) {$m$};
    \node[font=\scriptsize, simRust] at (0.0,-1.0) {$n$};
    \node[font=\scriptsize\ttfamily, simInk] at (0.0,-1.5) {Tensor[m, n]};
    \node[font=\scriptsize\itshape, simGray, text width=26mm, align=center]
      at (0.0,-2.05) {a dense matrix\\(rank 2)};
  \end{scope}
  % --- Tensor[H, W, 3]: stacked grids ---
  \begin{scope}[shift={(3.2,0)}]
    \def\dx{0.17}\def\dy{0.11}
    \foreach \k/\c in {0/simBgGray,1/simBgViolet,2/simBgRust}{
      \begin{scope}[shift={(\k*\dx,\k*\dy)}]
        \fill[\c, draw=simGray] (-0.8,-0.6) rectangle (0.8,0.6);
        \draw[simGray!55] (-0.8,-0.2)--(0.8,-0.2);
        \draw[simGray!55] (-0.8,0.2)--(0.8,0.2);
        \draw[simGray!55] (-0.27,-0.6)--(-0.27,0.6);
        \draw[simGray!55] (0.27,-0.6)--(0.27,0.6);
      \end{scope}}
    \node[font=\scriptsize, simInk] at (-1.15,0.0) {$H$};
    \node[font=\scriptsize, simInk] at (0.0,-0.95) {$W$};
    \node[font=\scriptsize, simRust] at (1.2,0.85) {$3$};
    \node[font=\scriptsize\ttfamily, simInk] at (0.15,-1.5) {Tensor[H, W, 3]};
    \node[font=\scriptsize\itshape, simGray, text width=28mm, align=center]
      at (0.15,-2.05) {a 3-vector per grid cell\\(rank 3)};
  \end{scope}
\end{tikzpicture}
\caption[A tensor as a labeled box of axes]{Three tensors, drawn as labeled
boxes of their axes. \emph{Left:} the global DoF vector
\lstinline[style=l4style]|Tensor[N]|, a single axis of length $N$---the shape
every solver vector wears. \emph{Center:} a small dense matrix
\lstinline[style=l4style]|Tensor[m, n]|, two axes. \emph{Right:} a field storing
a 3-vector at each cell of an $H\times W$ grid,
\lstinline[style=l4style]|Tensor[H, W, 3]|, three axes with the last pinned to
size $3$. The \emph{rank} is the number of axes; the bracketed list is the
\emph{shape}.}
\label{fig:t-box}
\end{figure}

\subsection{Real and complex scalars}
\label{sec:t-scalars}

Every tensor's entries are scalars of a single type, and our calculus admits
two: \emph{real} ($\Real$) and \emph{complex} ($\Complex$). The distinction is
physical, not cosmetic. The electrostatic and magnetostatic analyses
(\cref{ch:targets}) solve over $\Real$: a potential or vector-potential DoF is a
real number. The driven and eigenmode analyses are
\emph{time-harmonic}---a field oscillating as $e^{i\omega t}$ is most naturally a
\emph{complex} amplitude, so their DoF vectors are
\lstinline[style=l4style]|Tensor[N]| over $\Complex$. The transient analysis
marches real fields in time. A tensor's scalar type travels with it just as its
shape does, and operations must agree on it: you cannot silently add a real
tensor to a complex one. The natural question---when \emph{may} a real tensor be
treated as complex (the \emph{promotion} question)---we raise in
\cref{sec:t-promotion} and answer with the rest of the shape algebra in
\cref{ch:shapes}.

\section{Shapes as type-level contracts}
\label{sec:t-shapes}

A tensor's shape is a small piece of data---a bracketed list of axes---but it
plays the role of a \emph{type}. It is declared at construction, it travels with
the value, and it is \emph{checked at every function boundary}. A function that
expects a \lstinline[style=l4style]|Tensor[n]| and is handed a
\lstinline[style=l4style]|Tensor[m, n]| is as wrong as a function that expects an
integer and is handed a string---and just as with types, catching the mismatch at
the boundary, before any arithmetic runs, is what makes the calculus safe to
compose.

\subsection{The bracket notation}
\label{sec:t-bracket}

We write a shape as a comma-separated list of axes inside square brackets,
following the value's element type \code{Tensor}:

\begin{itemize}[leftmargin=1.4em, nosep]
\item \lstinline[style=l4style]|Tensor[N]| --- rank 1, a single axis of length
$N$. The flat DoF vector.
\item \lstinline[style=l4style]|Tensor[m, n]| --- rank 2, a small dense matrix
of $m$ rows and $n$ columns.
\item \lstinline[style=l4style]|Tensor[H, W, 3]| --- rank 3, with \emph{named}
axes $H$ and $W$ and a last axis \emph{pinned} to the literal size $3$ (a
3-vector per cell).
\end{itemize}

\noindent Two things in this notation are load-bearing. First, an axis may be a
\emph{name} (like $N$, $H$, $W$, $m$, $n$) or a \emph{literal} (like the $3$): a
name stands for a size we do not fix in the signature but require to
\emph{match} wherever the same name appears, while a literal pins the axis to an
exact extent. Second---and this is the point of giving axes \emph{names} rather
than just sizes---a shape records the \emph{meaning} of each axis, not only its
length. In \lstinline[style=l4style]|Tensor[E, P, C]| (\cref{ch:matrixfree}) the
axes are not interchangeable: $E$ counts elements, $P$ quadrature points, $C$
components, and a function that contracts the wrong one is wrong even if the
lengths happen to coincide. The name is how we keep them straight.

\begin{keyidea}
A shape is a \emph{type}. It states the rank, the size of each axis, and---through
named axes---the \emph{meaning} of each axis. It is part of the value, travels
with it, and is checked at every function boundary. Reading a function's shape
signature tells you what it consumes and produces before you read a line of its
body; a shape mismatch is a type error, caught at the boundary, not a wrong
number discovered later.
\end{keyidea}

\begin{notationbox}
When two operands must have the \emph{same whole shape of unspecified rank}, the
source calculus writes a single \emph{shape variable}---e.g.\
\lstinline[style=l4style]|axpy :: Scalar -> Tensor[$S] -> Tensor[$S] -> Tensor[$S]|---%
where \lstinline[style=l4style]|$S| names one common shape shared by all three.
This rank-agnostic ``same shape'' notation, and the richer \emph{named shape
groups} that let two tensors agree on \emph{part} of their shape, are the subject
of \cref{ch:shapes}. In this chapter we keep to concrete axes
(\lstinline[style=l4style]|Tensor[N]|, \lstinline[style=l4style]|Tensor[m, n]|)
and read \lstinline[style=l4style]|$S| informally as ``the same shape on both
sides.''
\end{notationbox}

\section{The primitive operations and their shape contracts}
\label{sec:t-primitives}

Here is where the linear algebra you already know meets typed tensors. Our
calculus is built from a small set of \emph{primitive} operations---the verbs of
\cref{ch:bigidea}'s first question. Each is a pure function, and each carries a
\emph{shape contract}: a signature that says exactly which shapes it consumes and
produces, and how their axes must line up. We give the six that recur most, each
as a signature and a shape-matching picture; their full algebra (laws, identities,
the operators that fold them together) is \cref{ch:operators}.

\begin{lstlisting}[style=l4style]
axpy    :: Scalar -> Tensor[$S] -> Tensor[$S] -> Tensor[$S]   -- a*x + y
dot     :: Tensor[$S] -> Tensor[$S] -> Scalar                 -- inner product
matvec  :: Tensor[m, n] -> Tensor[n] -> Tensor[m]             -- A . v
outer   :: Tensor[m] -> Tensor[n] -> Tensor[m, n]             -- v (X) w
reduce  :: Tensor[..., k] -> Tensor[...]                      -- collapse one axis
broadcast :: Tensor[$S_a] -> (target: Shape) -> Tensor[target]-- stretch to fit
\end{lstlisting}

\noindent Read each signature left to right: the argument shapes, then the
arrow, then the result shape. The two that share a shape (\code{axpy},
\code{dot}) demand their operands be \emph{identical} in shape; the two that
\emph{change} shape (\code{matvec}, \code{outer}) state the relation between input
and output axes explicitly; and the two that move \emph{between} ranks
(\code{reduce}, \code{broadcast}) add or drop an axis. \Cref{fig:t-shapematch}
draws the central one, \code{matvec}, as a shape-matching diagram.

\paragraph{\code{axpy} and \code{dot} --- same-shape contracts.}
\code{axpy} computes $a\vx+\vy$: both vectors and the result are one common shape
\lstinline[style=l4style]|$S| (in practice \lstinline[style=l4style]|Tensor[N]|),
and the scalar $a$ scales every entry. \code{dot} consumes two tensors of the
same shape and returns a single \code{Scalar}---it collapses \emph{all} axes by
summation. Over $\Complex$, \code{dot} conjugates its first argument (the
Hermitian inner product $\vx^{\mathsf H}\vy$); over $\Real$ this is the ordinary
$\sum_i x_i y_i$. The contract these share is the simplest and the most common:
\emph{the shapes must match exactly}, and a mismatch is rejected at the boundary.

\paragraph{\code{matvec} --- the contraction contract.}
\code{matvec} applies a small dense matrix to a vector. Its signature
\lstinline[style=l4style]|matvec :: Tensor[m, n] -> Tensor[n] -> Tensor[m]| says
the matrix's \emph{second} axis ($n$, its columns) must match the vector's single
axis ($n$), they are summed away in the contraction, and what survives is the
matrix's \emph{first} axis ($m$, its rows). This is exactly the rule from a first
linear-algebra course---``columns of $A$ must match the length of $\vv$, and the
product has as many entries as $A$ has rows''---now written as a shape contract
the calculus checks. \Cref{fig:t-shapematch} is the picture: the inner $n$'s
meet and cancel; the outer $m$ is the answer's shape.

\begin{figure}[t]
\centering
\begin{tikzpicture}[scale=1.0]
  % matrix [m, n]
  \node[draw=simGreen, fill=simBgGreen, thick, minimum width=20mm,
        minimum height=15mm] (A) at (0,0) {};
  \node[font=\scriptsize, simInk] at ($(A.west)+(-0.45,0)$) {$m$};
  \node[font=\scriptsize\bfseries, simRust] at ($(A.south)+(0,-0.45)$) {$n$};
  \node[font=\scriptsize\ttfamily, simInk, above=1mm of A] {Tensor[m, n]};
  % vector [n]
  \node[draw=simBlue, fill=simBgBlue, thick, minimum width=6mm,
        minimum height=15mm, right=12mm of A] (v) {};
  \node[font=\scriptsize\bfseries, simRust] at ($(v.south)+(0,-0.45)$) {$n$};
  \node[font=\scriptsize\ttfamily, simInk, above=1mm of v] {Tensor[n]};
  % equals / arrow
  \node[font=\Large] at ($(v.east)+(0.9,0)$) {$\Rightarrow$};
  % result [m]
  \node[draw=simRust, fill=simBgRust, thick, minimum width=6mm,
        minimum height=15mm, right=18mm of v] (r) {};
  \node[font=\scriptsize, simInk] at ($(r.west)+(-0.45,0)$) {$m$};
  \node[font=\scriptsize\ttfamily, simInk, above=1mm of r] {Tensor[m]};
  % the matching n's: a brace tying A's n and v's n
  \draw[decorate, decoration={brace, amplitude=4pt, mirror}, simRust]
    ($(A.south east)+(0,-0.75)$) -- ($(v.south west)+(0,-0.75)$)
    node[midway, below=4pt, font=\scriptsize, simRust]
    {inner $n$ must match, then cancels};
  % the surviving m
  \draw[-{Stealth[length=1.8mm]}, simGray, densely dashed]
    (A.west) to[out=170,in=170,looseness=1.3] (r.west);
  \node[font=\scriptsize\itshape, simGray] at (1.6,1.55)
    {outer $m$ survives as the result shape};
\end{tikzpicture}
\caption[Shape matching at a \code{matvec} boundary]{The shape contract of
\lstinline[style=l4style]|matvec :: Tensor[m, n] -> Tensor[n] -> Tensor[m]|, drawn
as axis matching. The matrix's column axis $n$ and the vector's single axis $n$
must agree (rust brace); the contraction sums over them and they vanish. The
matrix's row axis $m$ is untouched and becomes the shape of the result. A vector
of any length other than $n$ is rejected at the boundary, before a single
multiply-add runs.}
\label{fig:t-shapematch}
\end{figure}

\paragraph{\code{outer} --- the rank-raising contract.}
\code{outer} is \code{matvec}'s opposite: it takes two vectors and
\emph{builds} a matrix, $\vv\vw^{\mathsf T}$, whose $(i,j)$ entry is $v_i w_j$.
Its signature \lstinline[style=l4style]|outer :: Tensor[m] -> Tensor[n] -> Tensor[m, n]|
adds an axis rather than removing one: no axis is required to match, and the
result's rank is the \emph{sum} of the inputs'.

\paragraph{\code{reduce} and \code{broadcast} --- moving between ranks.}
\code{reduce} collapses one axis by an associative combiner (a sum, a max, a
product), turning a \lstinline[style=l4style]|Tensor[..., k]| into a
\lstinline[style=l4style]|Tensor[...]|: the labeled axis is summed (or
maxed, $\dots$) away and the rest of the shape survives. \code{dot} is the
special case that reduces \emph{every} axis to a scalar. \code{broadcast} runs the
other way: it \emph{stretches} a tensor to a larger target shape by repeating it
along new or length-1 axes---the mechanism by which a single scalar weight
multiplies a whole field, or a per-row factor scales every column. Its exact
matching rule (which shapes may be broadcast to which) is shape \emph{algebra},
and we defer it to \cref{ch:shapes}; here we note only that it is the
rank-\emph{raising} counterpart of \code{reduce}.

\begin{notationbox}
These six are not the whole vocabulary---\cref{ch:operators} adds the named
verbs (\code{nrm2}, \code{scal}, the operator-apply \code{apply\_linop}) and the
combinators that fold them together. But they are the \emph{primitives} whose
shape contracts every later operation is built from, and the discipline they
establish---\emph{a signature is a contract, checked at the boundary}---is the one
that makes the calculus composable.
\end{notationbox}

\subsection{Worked example: shape-checking a small 1-D solve}
\label{sec:t-wex-shapes}

Nothing fixes the contracts in the mind like running a real vector through them
and watching the shapes line up. We take the discretized 1-D field of Part~II---a
\lstinline[style=l4style]|Tensor[N]| of nodal values---and exercise \code{axpy},
\code{dot}, and a \code{matvec} against the tridiagonal stiffness matrix, checking
the shape at every boundary.

\begin{workedexample}{Shape-checked \code{axpy}, \code{dot}, \code{matvec} on a 1-D field}{t-shapes}
\textbf{Setup.} A 1-D mesh with $N=3$ interior nodes carries a discretized field
$\vx=(x_0,x_1,x_2)=(1,2,3)$ and a second field $\vy=(y_0,y_1,y_2)=(4,5,6)$, both
\lstinline[style=l4style]|Tensor[N]| with $N=3$ over $\Real$. The 1-D stiffness
operator (the discrete second difference, $-\partial_{xx}$ on a unit-spacing
mesh) is the tridiagonal matrix
\[
  \mat{K} \;=\;
  \begin{pmatrix} 2 & -1 & 0 \\ -1 & 2 & -1 \\ 0 & -1 & 2 \end{pmatrix},
  \qquad \mat{K} : \texttt{Tensor[m, n]},\ m=n=3 .
\]

\tcblower
\textbf{\code{axpy} (same-shape contract).} Compute $\vy' = a\vx + \vy$ with
$a=2$. The contract
\lstinline[style=l4style]|axpy :: Scalar -> Tensor[$S] -> Tensor[$S] -> Tensor[$S]|
demands both vectors share a shape: here \lstinline[style=l4style]|Tensor[3]| and
\lstinline[style=l4style]|Tensor[3]| agree, so the call type-checks, and
\[
  \vy' = 2\cdot(1,2,3) + (4,5,6) = (6,9,12) \;:\; \texttt{Tensor[3]} .
\]
The result wears the same shape as the inputs---rank and length are preserved.

\textbf{\code{dot} (reduce-to-scalar contract).} Compute $r=\vx\cdot\vy$. The
contract \lstinline[style=l4style]|dot :: Tensor[$S] -> Tensor[$S] -> Scalar|
again demands matching shapes (both \lstinline[style=l4style]|Tensor[3]|) and
collapses \emph{all} axes:
\[
  r = 1{\cdot}4 + 2{\cdot}5 + 3{\cdot}6 = 32 \;:\; \texttt{Scalar} .
\]
The shape \lstinline[style=l4style]|Tensor[3]| has gone to rank-0---the axis was
summed away.

\textbf{\code{matvec} (contraction contract).} Compute $\vw = \mat{K}\vx$. The
contract \lstinline[style=l4style]|matvec :: Tensor[m, n] -> Tensor[n] -> Tensor[m]|
checks that $\mat{K}$'s column axis $n=3$ matches $\vx$'s single axis $3$---it
does---then sums over them, leaving the row axis $m=3$:
\[
  \mat{K}\vx =
  \begin{pmatrix} 2{\cdot}1-1{\cdot}2 \\ -1{\cdot}1+2{\cdot}2-1{\cdot}3 \\ -1{\cdot}2+2{\cdot}3 \end{pmatrix}
  = (0,\,0,\,4) \;:\; \texttt{Tensor[3]} .
\]
The inner $3$'s cancelled (\cref{fig:t-shapematch}); the surviving $m=3$ is the
result shape, again a \lstinline[style=l4style]|Tensor[3]| field.

\textbf{The shapes line up.} Every call's boundary check passed because the axes
agreed: \code{axpy} and \code{dot} on two \lstinline[style=l4style]|Tensor[3]|
fields, \code{matvec} with the matrix's column count matching the vector's length.
Hand \code{matvec} a length-$2$ vector instead and the check fails immediately
($n=3\neq 2$), before any multiply-add---the type error caught at the boundary,
exactly as a wrong-typed argument would be.
\end{workedexample}

\section{Immutability makes aliasing harmless}
\label{sec:t-aliasing}

\Cref{ch:bigidea} argued that we describe computations as pure functions over
immutable values, and \cref{ch:records} carried the discipline into records. In
tensor terms the immutability has a concrete, daily payoff worth stating on its
own: \emph{aliasing is harmless}.

When two names refer to the same tensor, we say they \emph{alias} it. In a
mutable world, aliasing is a hazard---if \code{a} and \code{b} name the same
array and someone writes through \code{a}, then \code{b}'s value changes out from
under whoever was reading it, and the bug is among the nastiest in numerical
code. The whole apparatus of ``defensive copies'' exists to guard against exactly
this. But our tensors are immutable: \emph{no one can write through any name},
because there is no write at all---an operation produces a fresh tensor and the
old one stands. So two names pointing at one tensor can never interfere. Aliasing
costs nothing and risks nothing; we share freely and copy never.

\begin{figure}[t]
\centering
\begin{tikzpicture}[scale=1.0]
  % the one immutable tensor
  \node[draw=simBlue, fill=simBgBlue, thick, rounded corners=2pt,
        minimum width=26mm, minimum height=12mm] (t) at (0,0)
    {\footnotesize immutable tensor};
  \node[font=\scriptsize\ttfamily, simInk, below=1mm of t] {Tensor[N]};
  % two names pointing at it
  \node[font=\ttfamily\small, simRust] (na) at (-3.3,1.4) {x};
  \node[font=\ttfamily\small, simRust] (nb) at (-3.3,-1.4) {y};
  \draw[-{Stealth[length=2mm]}, simRust] (na) -- (t.north west);
  \draw[-{Stealth[length=2mm]}, simRust] (nb) -- (t.south west);
  \node[font=\scriptsize\itshape, simGray, align=left, text width=30mm]
    at (-3.6,0) {two names,\\one shared value};
  % the "would-be write" that never happens
  \node[font=\scriptsize, simGray, align=center, text width=34mm]
    at (4.0,0.7) {a ``modification'' of \texttt{x}\dots};
  \node[draw=simGreen, fill=simBgGreen, thick, rounded corners=2pt,
        minimum width=22mm, minimum height=10mm] (t2) at (4.0,-0.8)
    {\footnotesize \emph{new} tensor};
  \draw[-{Stealth[length=2mm]}, simGreen] (t.east) -- node[above, font=\scriptsize]
    {\texttt{f}} (t2.west);
  \node[font=\scriptsize\itshape, simGray, below=1mm of t2, text width=34mm,
        align=center] {\dots produces a fresh value;\\ \texttt{y} still sees the original};
\end{tikzpicture}
\caption[Aliasing is safe under immutability]{Two names, \texttt{x} and
\texttt{y}, share a single immutable tensor (blue). Because the tensor cannot be
written through any name, the sharing is invisible and free---no defensive copy
is needed. A function \texttt{f} that ``modifies'' \texttt{x} does not touch the
shared value at all; it \emph{returns a new tensor} (green), and \texttt{y}
continues to see the original unchanged. Mutation hazards simply cannot arise.}
\label{fig:t-aliasing}
\end{figure}

\begin{pitfall}
``No copy needed'' is a statement about \emph{meaning}, exactly as in
\cref{ch:bigidea}. It does \emph{not} mean an implementation literally keeps
every tensor forever: once no live name refers to a tensor, its storage may be
reclaimed---or, better, \emph{reused} in place for the next result, recovering the
mutation an experienced numericist would have written by hand. Immutability frees
the implementation to alias and to reuse \emph{because} it has proved the sharing
can never be observed. The discipline that makes this precise---which tensor's
storage may be reused for which result---is the state monad of \cref{ch:monad}.
\end{pitfall}

\section{The tensor-field lift}
\label{sec:t-lift}

We come to the chapter's central modeling move, and the true bridge from Part~II
to Part~III. Part~II's assembly chapters spoke the language of \emph{loops over
arrays}: ``for each degree of freedom $i$, do this to entry $i$.'' Part~III speaks
the language of \emph{operations on whole fields}: ``do this to the field.'' The
\emph{tensor-field lift} is the move that turns the first into the second.

\subsection{From a per-DoF loop to a field equation}
\label{sec:t-lift-what}

Take the \code{axpy} update once more. As a loop over the $N$ degrees of freedom
it reads, entry by entry,
\begin{equation}
  y_i \;\leftarrow\; a\,x_i + y_i, \qquad i = 0,1,\dots,N-1 .
  \label{eq:t-axpy-loop}
\end{equation}
The loop is explicit: it names an index $i$, ranges it over the DoFs, and updates
one entry per pass. The \emph{lift} observes that nothing in
\cref{eq:t-axpy-loop} couples one index to another---entry $i$ is computed from
entry $i$ alone---so the whole loop is captured, with nothing lost, by a single
statement about the field as a vector in $\Real^N$ (or $\Complex^N$):
\begin{equation}
  \vy' \;=\; a\,\vx + \vy .
  \label{eq:t-axpy-lift}
\end{equation}
The index $i$ has \emph{vanished}. \Cref{eq:t-axpy-lift} is one equation about
two whole tensors \lstinline[style=l4style]|x, y :: Tensor[N]| and a scalar $a$;
the per-DoF loop is now \emph{implicit}, recoverable but no longer written. This
is exactly the pure \code{axpy} of \cref{ch:bigidea}, and the lift is \emph{why}
its signature can speak of whole tensors without ever mentioning $i$.

\begin{keyidea}
The \emph{tensor-field lift} replaces a per-degree-of-freedom (or per-element)
loop with a single statement about a \emph{whole tensor field}. The update
$y_i \leftarrow a x_i + y_i$ \emph{for all $i$} becomes the one field equation
$\vy' = a\vx + \vy$ in $\Real^N$; the loop over $i$ is left implicit. This is the
move that lets the calculus describe a solver as operations on fields rather than
as nested loops over arrays---the single most important translation between
Part~II's assembly view and Part~III's calculus view.
\end{keyidea}

\begin{figure}[t]
\centering
\begin{tikzpicture}[scale=1.0]
  % LEFT: the explicit per-DoF loop
  \begin{scope}
    \node[font=\bfseries\small, simRust] at (0,2.3) {per-DoF loop};
    \node[draw=simRust, fill=simBgRust, rounded corners=2pt, thick,
          minimum width=30mm, minimum height=22mm, align=left, font=\scriptsize\ttfamily,
          text=simInk] (loop) at (0,0.5)
      {for i in 0..N:\\[2pt]\ \ y[i] = a*x[i]\\[1pt]\ \ \ \ \ \ \ \ + y[i]};
    \node[font=\scriptsize\itshape, simGray, below=1.5mm of loop, text width=34mm,
          align=center] {$N$ separate updates,\\index $i$ explicit};
  \end{scope}
  % collapse arrow
  \node[font=\Large, simBlue] at (3.3,0.5) {$\rightsquigarrow$};
  \node[font=\scriptsize\itshape, simBlue, align=center, text width=22mm]
    at (3.3,1.5) {lift\\(collapse the loop)};
  % RIGHT: the single field equation
  \begin{scope}[shift={(6.8,0)}]
    \node[font=\bfseries\small, simBlue] at (0,2.3) {field equation};
    \node[draw=simBlue, fill=simBgBlue, rounded corners=2pt, thick,
          minimum width=30mm, minimum height=22mm, align=center, font=\normalsize,
          text=simInk] (eq) at (0,0.5) {$\vy' = a\,\vx + \vy$};
    \node[font=\scriptsize\itshape, simGray, below=1.5mm of eq, text width=36mm,
          align=center] {one statement on whole tensors;\\ index $i$ gone, implicit};
  \end{scope}
\end{tikzpicture}
\caption[The tensor-field lift]{The lift collapses a per-DoF loop into a single
whole-field equation. \emph{Left (rust):} the explicit loop runs $N$ separate
updates, naming and ranging the index $i$. \emph{Right (blue):} the lifted form
is one equation about the whole tensors $\vx,\vy\in\Real^N$; the loop over $i$ has
disappeared into the vector notation. The two compute the same thing whenever no
index depends on another's result---the condition examined in
\cref{sec:t-lift-valid}.}
\label{fig:t-lift}
\end{figure}

\subsection{When the lift is valid}
\label{sec:t-lift-valid}

The lift is not magic, and it is not always legal. \Cref{eq:t-axpy-lift} is
faithful to \cref{eq:t-axpy-loop} precisely because of a structural fact: the
update at index $i$ reads inputs only at index $i$, never at some other index $j$
whose value the loop is in the middle of changing. There are two ways a per-DoF
loop can earn its lift.

\begin{description}[leftmargin=1.4em, style=nextline]
\item[Per-DoF independence.] Each output entry depends only on input entries at
the \emph{same} index (or on read-only data). \code{axpy}, \code{scal} ($\vx
\mapsto a\vx$), and elementwise products are all of this kind, and they lift to
whole-field equations directly. The weight stage $D$ of the matrix-free operator
(\cref{ch:matrixfree})---a pointwise multiply against \code{geom\_data}---is the
canonical multi-axis example: ``embarrassingly parallel,'' exactly because every
index is independent, exactly the condition for a clean lift.
\item[The coupling \emph{is} the operator.] Some operations \emph{do} couple
indices---but the coupling is the defining content of the operation, not a
loop-carried accident. \code{matvec}, and the full operator application
\code{apply\_linop} ($\vy = \mat{A}\vx$) of \cref{ch:matrixfree}, mix every input
entry into every output entry; that is what a linear operator \emph{is}. Such an
operation lifts to a single field equation $\vy = \mat{A}\vx$ because the coupling
is captured \emph{by the operator itself}, not by the order in which a loop visits
indices. There is no loop-carried dependency to preserve; the dependency is the
matrix.
\end{description}

\subsection{When the lift fails: a sequential reduction}
\label{sec:t-lift-fails}

The lift fails exactly when a loop \emph{carries a dependency forward}: the update
at index $i$ reads a value the \emph{same loop} produced at an earlier index, so
the order of the loop is part of the answer. Two small examples make the failure
concrete.

A \emph{running sum} (prefix sum) over a tensor,
\begin{equation}
  s_i \;=\; s_{i-1} + x_i, \qquad s_0 = x_0,
  \label{eq:t-prefix}
\end{equation}
has no whole-field equation of the form ``$\vs = (\text{something})\,\vx$'' that
respects it as a single independent statement, because $s_i$ \emph{needs}
$s_{i-1}$---a value computed one step earlier in the same sweep. Try to write it
as ``do all indices at once'' and you cannot: index $i$ has nothing to read until
index $i-1$ is finished. The order is load-bearing.

The same obstruction, in the solver world, is a \emph{Gauss--Seidel} sweep, where
updating component $i$ uses the \emph{already-updated} components $j<i$ from the
\emph{current} sweep:
\begin{equation}
  x_i \;\leftarrow\; \frac{1}{a_{ii}}\Big(b_i - \sum_{j<i} a_{ij}\,x_j^{\text{new}}
       - \sum_{j>i} a_{ij}\,x_j^{\text{old}}\Big).
  \label{eq:t-gs}
\end{equation}
The term $x_j^{\text{new}}$ for $j<i$ is the loop reading its own fresh output:
the very loop-carried dependency the lift cannot absorb. \Cref{fig:t-liftfail}
contrasts the two situations side by side---the independent fan (lifts) against
the sequential chain (does not).

\begin{figure}[t]
\centering
\begin{tikzpicture}[scale=1.0]
  % LEFT: independent — lifts
  \begin{scope}
    \node[font=\bfseries\small, simBlue] at (0,2.1) {independent --- lifts};
    \foreach \i in {0,1,2,3}{
      \node[draw=simBlue, fill=simBgBlue, circle, minimum size=6mm,
            font=\scriptsize] (xi\i) at (\i*0.95,1.1) {$x_{\i}$};
      \node[draw=simBlue, fill=simBgBlue, circle, minimum size=6mm,
            font=\scriptsize] (yi\i) at (\i*0.95,-0.4) {$y'_{\i}$};
      \draw[-{Stealth[length=1.6mm]}, simBlue] (xi\i) -- (yi\i);
    }
    \node[font=\scriptsize\itshape, simGray, align=center, text width=40mm]
      at (1.4,-1.25) {each output reads its own index only\\
      $\Rightarrow$ collapses to $\vy'=a\vx+\vy$};
  \end{scope}
  % RIGHT: sequential — does not lift
  \begin{scope}[shift={(6.1,0)}]
    \node[font=\bfseries\small, simRust] at (0.4,2.1) {sequential --- does \emph{not} lift};
    \foreach \i in {0,1,2,3}{
      \node[draw=simRust, fill=simBgRust, circle, minimum size=6mm,
            font=\scriptsize] (si\i) at (\i*0.95,0.35) {$s_{\i}$};
    }
    \foreach \i/\j in {0/1,1/2,2/3}{
      \draw[-{Stealth[length=1.8mm]}, simRust, thick] (si\i) -- (si\j);}
    \node[font=\scriptsize\itshape, simGray, align=center, text width=42mm]
      at (1.4,-1.25) {$s_i$ needs $s_{i-1}$ --- a chain\\
      $\Rightarrow$ order is load-bearing, no field equation};
  \end{scope}
\end{tikzpicture}
\caption[When the lift fails]{The structural difference between a liftable and an
unliftable loop. \emph{Left (blue):} an \code{axpy}-style update fans
independently---each $y'_i$ reads only $x_i$ (and $y_i$), no arrow runs between
output indices---so the loop collapses to the field equation $\vy'=a\vx+\vy$.
\emph{Right (rust):} a running sum $s_i=s_{i-1}+x_i$ is a \emph{chain}---each
node depends on the one before---so the order of evaluation is part of the answer
and no single whole-field equation captures it. Sequential reductions of this
kind (Gauss--Seidel, triangular solves) are \emph{obstructions} to the lift, and
\cref{part:framework} records them as such rather than forcing a false collapse.}
\label{fig:t-liftfail}
\end{figure}

When the lift fails, we do \emph{not} pretend it succeeded. The honest record is
an \emph{obstruction}: a note that this loop carries a sequential dependency that
no whole-field statement captures, with the specific recurrence cited. Part~IV's
solver chapters meet several---the Gauss--Seidel and SSOR smoothers, the
triangular solves inside a direct factorization---and each is logged as a
sequential obstruction rather than lifted. (Forward reference: the iterative
solvers of \cref{part:framework} and the smoothers behind the multigrid
preconditioner are where these surface in force.) The lift is a powerful default,
not a universal law; knowing \emph{when} it fails is as important as knowing how
it works.

\subsection{Worked example: lifting a per-DoF update, and one that won't lift}
\label{sec:t-wex-lift}

We make the lift, and its failure, fully concrete on a tiny $N=4$ field: write a
per-DoF update first as an explicit loop, then as a single tensor-field equation,
and argue they agree---then exhibit one update where the argument breaks down.

\begin{workedexample}{The lift: a loop and a field equation that agree, and one that cannot}{t-lift}
\textbf{Setup.} Take $N=4$ and the fields $\vx=(1,2,3,4)$, $\vy=(10,20,30,40)$
as \lstinline[style=l4style]|Tensor[4]| over $\Real$, with scalar $a=3$.

\tcblower
\textbf{The update as an explicit loop.} The \code{axpy} update visits each index
in turn:
\begin{lstlisting}[style=l4style, language=]
for i in 0..4:
    yp[i] = a * x[i] + y[i]
-- i=0: 3*1+10 = 13     i=1: 3*2+20 = 26
-- i=2: 3*3+30 = 39     i=3: 3*4+40 = 52
\end{lstlisting}
producing $\vy'=(13,26,39,52)$. Four updates, the index $i$ named and ranged.

\textbf{The same update as a field equation.} The lift collapses the loop to one
statement on whole tensors,
\[
  \vy' \;=\; a\,\vx + \vy \;=\; 3\,(1,2,3,4) + (10,20,30,40) \;=\; (13,26,39,52),
\]
the very same four numbers. The index $i$ never appears.

\textbf{Why they agree.} Entry $i$ of the loop reads $x_i$ and $y_i$ and writes
$y'_i$---inputs and output \emph{all at the same index}, with no entry of $\vy'$
read before it is written. The updates are mutually independent, so doing them
``all at once'' as a vector equation reorders nothing that matters: this is the
per-DoF-independence condition of \cref{sec:t-lift-valid}, and the lift is exact.

\textbf{One that will not lift.} Now replace the body with a \emph{running sum}:
\begin{lstlisting}[style=l4style, language=]
s[0] = x[0]
for i in 1..4:
    s[i] = s[i-1] + x[i]      -- reads s[i-1]: this loop's own output
-- s = (1, 3, 6, 10)
\end{lstlisting}
Each $s_i$ reads $s_{i-1}$, a value \emph{this same loop produced one step
earlier}. There is no equation $\vs = a\vx + (\dots)$ that computes all four
entries independently and respects the recurrence: index $i$ has nothing to read
until index $i-1$ has finished. The order of evaluation is part of the answer, so
the lift \emph{fails}, and the honest record is a sequential
\emph{obstruction}---not a forced field equation
(\cref{fig:t-liftfail}). The same obstruction is the Gauss--Seidel sweep of
\cref{eq:t-gs}; \cref{part:framework} gives these their proper home in the fold
(\cref{ch:fold}).
\end{workedexample}

\section{Scalars as rank-0 tensors, and the promotion question}
\label{sec:t-scalars-rank0}

One loose end ties the chapter off and sets up the next. A single
scalar---the $a$ in \code{axpy}, the result of \code{dot}, a material
coefficient---is a tensor too: a \emph{rank-0} tensor, a box with \emph{no} axes
and exactly one entry. This is not a pedantic flourish. Treating a scalar as a
rank-0 tensor is what lets \code{broadcast} (\cref{sec:t-primitives}) stretch it
across a whole field uniformly: ``multiply the field by $a$'' is ``broadcast the
rank-0 $a$ to the field's shape, then multiply pointwise.'' The scalar and the
field meet through the same shape machinery as any two tensors.

\subsection{The promotion question, set up}
\label{sec:t-promotion}

That leaves the question \cref{sec:t-scalars} flagged: when may a \emph{real}
tensor be used where a \emph{complex} one is expected? The need is real and
concrete. A driven analysis (\cref{ch:targets}) works over $\Complex$, but its
mesh geometry, its material coefficients, and an initial real excitation are
naturally real; somewhere they must be \emph{promoted} into the complex
computation without a clumsy manual rewrite. Intuitively the promotion is the
obvious embedding $r \mapsto r + 0i$, and it should compose silently---adding a
real tensor to a complex one should ``just work,'' yielding a complex result.

But ``just work'' hides real questions. \emph{Which} direction promotes (real
into complex, never the reverse)? Does the result shape change (no---only the
scalar type)? When two operands disagree on \emph{both} shape and scalar type, in
which order are they reconciled? These are questions of \emph{shape and type
algebra}, and answering them precisely---alongside broadcasting's exact matching
rule and the named-shape-group notation for partial congruence---is the whole
business of the next chapter.

\begin{notationbox}
We have now met every piece the shape algebra operates on: tensors of a rank and
a scalar type, axes that are names or literals, the same-shape and contraction
contracts, rank-0 scalars, and the promotion of real into complex. What we have
\emph{not} done is give the rules by which shapes \emph{combine}---how broadcasting
decides two shapes are compatible, how a named shape group asserts partial
congruence, how scalar-type promotion composes with shape matching.
\Cref{ch:shapes} is exactly that algebra. This chapter fixed the vocabulary; the
next makes it compute.
\end{notationbox}

\summarysection
A \emph{tensor} in our calculus is an immutable, typed, multi-axis array of
scalars (real or complex); its \emph{rank} is its number of axes and its
\emph{shape} is the bracketed list of axis sizes and meanings. Every quantity
Part~II built is one: a discretized field is the rank-1 global DoF vector
\lstinline[style=l4style]|Tensor[N]|, and an element-local quantity
(\cref{ch:matrixfree}) is a small multi-axis tensor like
\lstinline[style=l4style]|Tensor[E, P, C]|. A shape is a \emph{type}: it travels
with the value and is checked at every function boundary, so a mismatch is caught
before any arithmetic runs. The primitive operations carry \emph{shape
contracts}---\code{axpy} and \code{dot} demand identical shapes, \code{matvec}
($[m,n]\cdot[n]\to[m]$) and \code{outer} relate input and output axes,
\code{reduce} and \code{broadcast} move between ranks. Because tensors are
immutable, \emph{aliasing is harmless}: two names may share one tensor with no
risk and no copy. The chapter's central move is the \emph{tensor-field lift}: a
per-DoF loop $y_i \leftarrow a x_i + y_i$ becomes the single field equation
$\vy'=a\vx+\vy$, the index left implicit. The lift is valid when each index is
independent, or when the coupling \emph{is} the operator (\code{apply\_linop}); it
\emph{fails} for a sequential reduction (a running sum, a Gauss--Seidel sweep),
which we record as an obstruction rather than force. Scalars are rank-0 tensors,
and the promotion of real into complex is set up here and made precise---with the
full shape algebra---in \cref{ch:shapes}.

\furtherreading
The matrix--vector and inner-product contracts of \cref{sec:t-primitives} are the
elementary operations of numerical linear algebra; Trefethen and
Bau~\citep{trefethenbau1997} give the cleanest conceptual account of why the
\emph{action} $\vv\mapsto\mat{A}\vv$, not the entries, is what algorithms consume
(the thread we pick up in \cref{ch:matrixfree}), and Golub and Van
Loan~\citep{golubvanloan2013} are the standard comprehensive reference for the
shapes and costs of these operations. The named-axis, shape-as-type discipline
sketched here is the same one modern tensor frameworks enforce at array
boundaries; we develop its algebra from scratch in \cref{ch:shapes} and apply it
to the calculus's verbs in \cref{ch:operators}, so no outside background is
assumed. The tensor-field lift and its sequential-obstruction failure mode return
throughout \cref{part:framework}, where the fold (\cref{ch:fold}) gives the
unliftable chains a home of their own.

\exercisessection
\begin{enumerate}[nosep]
  \item For each of \lstinline[style=l4style]|Tensor[N]|,
  \lstinline[style=l4style]|Tensor[m, n]|, and
  \lstinline[style=l4style]|Tensor[E, P, C]|, state the rank and say in one
  phrase what each axis \emph{counts}. Why is naming the axes (rather than only
  giving their lengths) load-bearing for the last one?
  \item The contract for \code{matvec} is
  \lstinline[style=l4style]|Tensor[m, n] -> Tensor[n] -> Tensor[m]|. A colleague
  hands the function a matrix of shape \lstinline[style=l4style]|Tensor[m, n]| and
  a vector of shape \lstinline[style=l4style]|Tensor[m]|. Exactly where does the
  shape check fail, and why is failing \emph{at the boundary} better than
  producing a wrong number?
  \item Write the \code{scal} operation $\vx \mapsto a\vx$ as (a) an explicit
  per-DoF loop and (b) a single field equation. Argue that the two compute the
  same thing, naming the structural property of the loop that licenses the lift.
  \item Give the shape contract (a signature
  \lstinline[style=l4style]|Tensor[...] -> ... -> Tensor[...]|) for an
  \emph{elementwise product} of two same-shape tensors, and for \code{outer}.
  Which one raises the rank, which preserves it, and why?
  \item Consider the prefix-sum recurrence $s_i = s_{i-1} + x_i$. Explain
  precisely why it has no whole-field equation of the form $\vs = \mat{A}\vx$ that
  honors it as a single independent statement. Contrast this with \code{matvec},
  which \emph{does} couple every index yet still lifts---what is different about
  the kind of coupling?
  \item Two names \code{x} and \code{y} alias one immutable
  \lstinline[style=l4style]|Tensor[N]|. A function returns a fresh tensor
  ``modifying'' \code{x}. State what \code{y} sees afterward, and explain why an
  implementation is nonetheless free to reuse \code{x}'s storage for the new
  result. What condition must hold for that reuse to be safe?
  \item A driven analysis (\cref{ch:targets}) computes over $\Complex$ but reads
  a real material coefficient. Describe the promotion $r \mapsto r + 0i$ as a map
  on scalar type, and say what it does (and does not) change about a tensor's
  shape. Why must promotion run real-into-complex and never the reverse?
  \item \emph{(Looking ahead.)} The weight stage $D$ of the matrix-free operator
  (\cref{ch:matrixfree}) is a pointwise multiply on a
  \lstinline[style=l4style]|Tensor[E, P, C]|. Argue that it lifts cleanly to a
  whole-field equation, identifying which of the two validity conditions of
  \cref{sec:t-lift-valid} it satisfies. Then explain why the \emph{scatter}
  stage $G^{\mathsf T}$, which sums shared-DoF contributions, is a different
  matter.
\end{enumerate}
